\documentclass[10pt,a4paper]{amsart}
\usepackage[margin=19mm]{geometry}
\usepackage{amsmath,amssymb,mathtools}
\usepackage[scr=rsfs,cal=euler]{mathalfa}
\usepackage{xfrac}
\usepackage[unicode,hidelinks,bookmarks=false]{hyperref}
\newcommand{\ie}{i.e.\ }
\newcommand{\cf}{cf.\ }
\newcommand{\resp}{respectively\ }
\newcommand{\Subsection}{\subsection}
\providecommand{\nobreakdash}{\nobreak\hbox{-}}
\setlength{\emergencystretch}{2em}
\multlinegap0pt
\usepackage{mathtools}
\usepackage{amssymb}
\usepackage[shortlabels]{enumitem}
\usepackage{tikz}
\newtheorem{lemma}{Lemma}
\newcommand{\C}{\mathcal{C}} % C_b(S\times S)
\newcommand{\R}{\mathbb{R}} % Real numbers
\title{A Distributionally Robust Framework for Learned Reconstructions in Inverse Problems}
\author{Floor van Maarschalkerwaart, Subhadip Mukherjee, Christoph Brune, Marcello Carioni}
\date{Source: arXiv:2606.30230v1 (CC BY 4.0)}
\begin{document}
\maketitle

\noindent\emph{Source and licence.} A Distributionally Robust Framework for Learned Reconstructions in Inverse Problems. Version arXiv:2606.30230v1 (CC BY 4.0), distributed by arXiv under the Creative Commons Attribution 4.0 International licence, http://creativecommons.org/licenses/by/4.0/, as recorded in the arXiv licence metadata for this exact version and shown on the abstract page https://arxiv.org/abs/2606.30230v1. Full licence notice: \texttt{licenses/arxiv-2606.30230-CC-BY-4.0.txt}.\par\smallskip

\noindent\textit{Editorial scope.} Counted unit: the article's lemma lemma:fenchel-requirements (source0.tex lines 1218--1220). The article's complete original proof of it is delivered with it (source0.tex lines 1221--1286, one \texttt{proof} environment of 66 lines). No mathematics is written, repaired or re-derived: the statement and the proof are the article's own lines, in the article's own notation, and no retained line is substituted or re-flowed.  Retained uncounted as the source-local context the proof needs: lines 1164--1167.  Nothing was omitted from any retained span, so the package carries no omission record.  No retained line contains a citation, so the wrapper carries no bibliography and no bibliography entry.  No retained line contains a figure environment, an image-inclusion command or drawing code, so no image is delivered and the package declares no asset dependency.  Editorial formatting only, in addition: an AMS \texttt{amsart} preamble that restates the article's own declarations and macros used by the retained lines, namely the environments \texttt{lemma} on separate counters, together with the standard packages those lines need.  The retained environments print in this document as Lemma 1 in order of appearance, whereas the article prints the same items with the numbering of its own full document.  The complete native source of the article as distributed in the arXiv e-print for this exact version is bundled unchanged as \texttt{sources/204001/source0.tex}.
\begin{align}\label{eq:B}
    B:= \bigg\{ g \in \C: \int_S\int_S &g(s,r)d\pi_{S|S=s}(r)d\mu^*(s) \\
    &\geq \int_S\int_S\ell(r,G) d\pi_{S|S=s}(r) d\mu^*(s), \ \forall \pi_{S|S}: S\rightarrow K_s \bigg\}\nonumber.
\end{align}

\begin{lemma}\label{lemma:fenchel-requirements}
    Let $S$ be a compact Polish space, and assume $c:S\times S\to\R_+$ is non-negative lower semi-continuous such that $c(s,r)=0$ if and only if $s=r$ and $\ell$ is continuous. Define  $B$ as in \eqref{eq:B}. Then, $B$ has non-empty interior.
\end{lemma}

\begin{proof}
Since $S$ is compact and
$\ell$ is continuous, there exists
\begin{align*}
M:=\sup_{r\in S}\ell(r,G)<+\infty .
\end{align*}
Define $k(s,r):=M+1$ for $ (s,r)\in S\times S$.
Then \(k\in C(S\times S)\). Moreover, for every $\pi_{S|S} :S \rightarrow K_s$ it holds that
\begin{align*}
\int_{S\times S} k(s,r)\,d\pi_{S|S=s}(r)d\mu^*(s)
=
M+1
\geq
\int_{S\times S}\ell(r,G)\,d\pi_{S|S=s}(r)d\mu^*(s),
\end{align*}
and therefore $k\in B$.
We claim that $k\in \operatorname{int}(B)$. Indeed, let
$\eta\in C(S\times S)$ be such that
\begin{align*}
\|\eta-k\|_\infty<\frac12 .
\end{align*}
Then, for every \((s,r)\in S\times S\),
\begin{align*}
\eta(s,r)
\geq
k(s,r)-\frac12
=
M+\frac12
\geq
\ell(r),
\end{align*}
so that $\eta \in B$ and thus $k\in\operatorname{int}(B)$.

%    Denote by $ri(X)$ the relative interior of a set $A$. We will show there exists an $k \in A \cap B$ such that $k \in ri(A) \cap ri(B)$. We take $k(s,r) := c(s,r) + \sup_{s \in S} \ell (s)$, where $\sup_{z \in S} \ell (z) < \infty$, since $\ell$ is upper semi-continuous and $S$ is compact. Choosing
%    \begin{align}
%       h(s,r) =  \sup_{z \in S} \ell (z)  \quad \text{ and } \quad \lambda = 1
%    \end{align}
%    we find $ h(s,r)+\lambda c(s,r)=k(s,r)$ and thus 
%    \begin{align}
%       \int_S\int_S k(s,r)d\pi_{S|S=s}(r)d\mu^*(s) = \int_S\int_S h(s,r) + \lambda c(s,r) d\pi_{S|S=s}(r) d\mu^*(s), \quad \forall \pi : S \rightarrow K
%    \end{align}
%    so $k \in A$. Since $k(s,r) \geq \ell (r)$ for every $r$, we also have
%    \begin{align}
%        \int_S\int_S k(s,r) d\pi_s(r)d\mu^*(s) \geq \int_S\int_S \ell(r)d\pi_s(r)d\mu^*(s), \quad \forall \pi : S \rightarrow K
%    \end{align}
%    and thus $k \in B$ and also $k \in A \cap B$. 
    
%    To show $k \in ri(A)$ we need only show that for every $g \in A$ there exists a $\tau > 1$ such that $\tau k +(1-\tau)g \in A$, and analogous for $B$ (since both $A$ and $B$ are convex). 
%    Let $\tau >1$ and $g \in A$ arbitrary, but different than $k$. Define $f(s,r):= \tau k(s,r) +(1-\tau)g(s,r).$ Then for every $\pi :S \rightarrow K$ it holds
%    \begin{align}
%        \int_S\int_S f(s,r)d\pi_s(r)d\mu^*(s) &= \int_S\int_S \tau(c(s,r) +\sup_{s \in S} \ell (s))\ d\pi_s(r)d\mu^*(s) + \int_S\int_S (1-\tau)g(s,r)\ d\pi_s(r)d\mu^*(s) \\
%        &= \int_S\int_S \tau(c(s,r) + \sup_{z \in S} \ell(z))\ d\pi_s(r)d\mu^*(s) \\
%        &\qquad+ \int_S\int_S (1-\tau)(h(s,r)+c(s,r))\ d\pi_s(r)d\mu^*(s) \\
%        &= \int_S\int_S [(1-\tau)h(s,r) + \tau \sup_{z \in S} \ell(z) +  c(s,r)]\ d\pi_s(r)d\mu^*(s) \\
%        &= \int_S\int_S [h_2(s,r) + \lambda_2 c(s,r)]\ d\pi_s(r)d\mu^*(s)
%    \end{align}
%    where we define $h_2(s,r) = (1-\tau)h(s,r) + \tau \sup_{z \in S} \ell(z) \in \mathcal{C}$ and $\lambda_2:= 1$ so that $f \in A$ and thus $h \in ri(A)$. 
%    Now let $g \in B$ arbitrary and again $\tau > 1$. Define similarly $f(s,r):= \tau k(s,r) + (1-\tau)g(s,r)$. Then for every $\pi : S \rightarrow K$ it holds
%    \begin{align}
%        \int_S\int_S f(s,r)\ d\pi_s(r)d\mu^*(s) &=  \int_S\int_S \tau(c(s,r)+\sup_{z \in S} \ell(z))\ d\pi_s(r)d\mu^*(s) + \int_S\int_S (1 - \tau)g(s,r)\ d\pi_s(r)d\mu^*(s) \\
%        &\geq \int_S\int_S \tau(c(s,r) + \ell(r))\ d\pi_s(r)d\mu^*(s) + \int_S\int_S (1-\tau) \ell(r)\ d\pi_s(r)d\mu^*(s) \\
%        &= \int_S\int_S \tau c(s,r) + \ell (r) \ d\pi_s(r)d\mu^*(s) \geq \int_S\int_S \ell(r) \ d\pi_s(r)d\mu^*(s)
%    \end{align}
%    so $f \in B$ and thus $k \in ri(B)$. We conclude that $ri(A) \cap ri(B)$ thus is non-empty. 

\end{proof}

\end{document}
